\documentclass[11pt,reqno]{amsart}

\usepackage{amsmath,amssymb,amsthm}
\usepackage[margin=1.1in]{geometry}
\usepackage{microtype}
\usepackage[colorlinks=true,linkcolor=blue,citecolor=blue]{hyperref}
\raggedbottom

\theoremstyle{plain}
\newtheorem{theorem}{Theorem}[section]
\newtheorem{proposition}[theorem]{Proposition}
\newtheorem{lemma}[theorem]{Lemma}
\newtheorem{corollary}[theorem]{Corollary}
\theoremstyle{definition}
\newtheorem{definition}[theorem]{Definition}
\newtheorem{remark}[theorem]{Remark}

\newcommand{\R}{\mathbb{R}}
\newcommand{\Z}{\mathbb{Z}}
\newcommand{\F}{\mathbb{F}}
\newcommand{\pill}{P}
\newcommand{\iss}{\looparrowright}

\begin{document}

\title[Exact action separation in the pillowcase]
{Exact action separation for a pretzel-tangle composition in the pillowcase}

\author{Bernd J. Wuebben}
\subjclass[2020]{53D12 (primary); 53D40, 57K18, 57K31 (secondary)}
\keywords{Pillowcase, traceless character variety, Lagrangian correspondence,
immersed curve, action filtration, rational tangle}
\date{August 15, 2026}

\begin{abstract}
Let $P$ be the smooth stratum of the traceless $\mathrm{SU}(2)$ character variety of the
four-punctured sphere.  We study the immersed curve obtained by composing the rational-tangle
arcs for $Q_{1/3}$ and $Q_{1/7}$ with the small-holonomy-perturbation correspondence associated
to a three-strand tangle.  After fixing an equivariant Liouville completion of $P$, we prove that
the truncated correspondence is exact for every sufficiently small positive perturbation and
that a half-period symmetry fixes the relative primitive on the two components of the output.

The singular composition consists of $21$ separated graph sheets joined across six fold circles.
Their twelve endpoints have nondegenerate hyperbolic smoothing germs.  From the resulting exact
path word we enumerate $58$ seam-overlap candidates.  Rational interval certificates for the
normal-displacement equations retain exactly $50$ transverse ordered self-intersection generators.
An area-residue calculation evaluates every limiting action as a rational multiple of $\pi^2$:
$42$ are negative and $8$ are positive, no action vanishes, and the smallest absolute limiting
action is $\pi^2/42$.  The census and all action signs persist for sufficiently small positive
perturbation.  The proof is accompanied by a pure-rational executable certificate.  No
holomorphic-quilt count, bounding cochain, or instanton-pillowcase comparison is asserted.
\end{abstract}

\maketitle

\section{Introduction}

The traceless $\mathrm{SU}(2)$ character variety of the four-punctured sphere is the pillowcase
\[
 P=(S^1\times S^1)/(\gamma,\theta)\sim(-\gamma,-\theta).
\]
Its smooth stratum $P^*$ is obtained by deleting the four orbifold points; we use the standard
Liouville completion near those punctures \cite[Appendix~A.1.1]{CHKK}.  Rational tangles give
proper Lagrangian arcs $A_r\subset P^*$, while a holonomy perturbation of the three-strand tangle
gives an immersed Lagrangian correspondence
\[
 C_t=C_{3,t}^{\mathrm{reg}}\iss(P_1^*)^-\times(P_2^*)^-\times P_3^*
\]
for small $t>0$ \cite{HeraldKirk,Smith}.  We consider the direct fiber product and its output immersion
\begin{equation}\label{eq:composition}
 \widetilde L_t=(A_{1/3}\times A_{1/7})
       \times_{P_1\times P_2} C_t,
 \qquad
 L_t=\pi_3(\widetilde L_t)\iss P_3^*.
\end{equation}
The pair of slopes $1/3$ and $1/7$ is the character-variety composition associated to the
pretzel-tangle presentation of $P(-2,3,7)$ used in the companion calculation \cite{WuebbenII}.

The purpose of this paper is finite and geometric.  It determines the exact immersed-curve action
data of \eqref{eq:composition}; it does not construct any moduli space of figure-eight bubbles.
This distinction matters because strip shrinking can produce Morsified trees of several
holomorphic components \cite{BottmanWehrheim}.  A list of candidate output generators is not a
count of those trees.

We fix the following terminology.  If a transverse self-intersection of $L_t$ has two preimages
$z_-$ and $z_+$, the orientation convention for immersed Floer theory selects one ordered pair
$x=(z_-,z_+)$ in degree one.  We call this ordered generator a \emph{root}.  If $f_t$ is the
primitive inherited from the exact fiber product, its action is
\begin{equation}\label{eq:root-action-intro}
 \mathcal A_t(x)=f_t(z_-)-f_t(z_+).
\end{equation}
The word ``root'' in this paper therefore means an ordered branch jump, not a zero of an
unspecified function.  The \emph{short deck class} of a root is the reduced torus-deck
displacement of the ordinary closing loop used in the area-residue calculation; for an odd seam
separation, the construction first doubles the path and then removes that artificial factor of
two.

\begin{theorem}[Exact action census]\label{thm:main}
There is an equivariant Liouville completion of the pillowcase and a number $t_0>0$ with the
following properties for every $0<t<t_0$.
\begin{enumerate}
\item The compactly truncated correspondence $C_t$ is exact.  The direct output $L_t$ inherits a
primitive whose relative normalization across its two open parameter arcs is fixed by a
half-period symmetry.
\item The curve $L_t$ has exactly $50$ transverse degree-one roots arising from the singular
composition: $28$ diagonal-circle roots and $22$ circle-circle roots.  They occur in $25$
half-period pairs.
\item The limiting actions as $t\to0^+$ are nonzero rational multiples of $\pi^2$.  Exactly $42$
roots have negative limiting action and $8$ have positive limiting action.  The minimum absolute
limiting action is $\pi^2/42$, and all $50$ action signs persist throughout $0<t<t_0$.
\end{enumerate}
The four positive half-period representatives have short deck classes and limiting actions
\[
\begin{array}{c|c|c}
\text{seam pair}&\text{short deck}&\displaystyle\lim_{t\to0^+}\mathcal A_t/\pi^2\\ \hline
(1,7)&(3,1)&3/7\\
(2,7)&(8,4)&5/14\\
(2,17)&(3,1)&3/14\\
(3,4)&(10,5)&5/42.
\end{array}
\]
\end{theorem}

Three ingredients make the statement exact.  First, the Liouville primitive is normalized
equivariantly and its periods are checked on the full singular fiber, including seam-born and
corner-link cycles.  Hamiltonian-composition flux then transports this vanishing to $t>0$.
Second, the singular output is recorded as an exact piecewise-linear word.  Its only failures of
the graph description occur at six explicitly determined stabilizer circles, whose endpoints
carry hyperbolic smoothing germs.  Third, all circle-circle intersections reduce to quadratic
equations over $\mathbb Q(2\cos(\pi/7))$.  Rational enclosures certify the physical roots and their
derivative signs.  The action of each surviving root is then a signed area minus a lattice-residue
term.

Section~\ref{sec:exactness} fixes the primitive and proves exactness.  Section~\ref{sec:geometry}
describes the six folds, the hyperbolic smoothing, and the $21$ separated graph sheets.
Section~\ref{sec:actions} proves the exact root census and action table.
Section~\ref{sec:scope} records the precise energy consequence and the boundary of the result.

\section{Exact primitives and half-period normalization}\label{sec:exactness}

Write $(\gamma,\theta)$ for the torus-cover coordinates on the pillowcase.  On the central
cylinder
\[
 U=(0,\pi)\times S^1\subset P^*,
\]
the Goldman form is $d\gamma\wedge d\theta$.  The choice of a global primitive on the completed
four-punctured sphere includes residue data at its four ends.

\begin{lemma}[Normalized pillowcase completion]\label{lem:normalized-primitive}
Let
\[
 \tau[\gamma,\theta]=[\gamma,\theta+\pi],
 \qquad
 \sigma[\gamma,\theta]=[\pi-\gamma,-\theta].
\]
There is a $\langle\tau,\sigma\rangle$-invariant corner truncation
$P^\delta\subset P^*$ and a $\langle\tau,\sigma\rangle$-invariant Liouville completion
$(\widehat P^\delta,\lambda)$ such that
\[
 d\lambda=d\gamma\wedge d\theta\quad\text{on }P^\delta,
 \qquad
 \lambda=(\gamma-\pi/2)\,d\theta
\]
near the circle $\{\gamma=\pi/2\}$.  The four oriented end residues of $\lambda$ are equal.
\end{lemma}

\begin{proof}
Choose the four deleted disks equivariantly.  Since the compact surface with boundary $P^\delta$
has zero second cohomology, the restricted Goldman form has a primitive $\lambda'$.  On an annular
neighborhood of $\{\gamma=\pi/2\}$, add a closed one-form to kill the period of
$\lambda'-(\gamma-\pi/2)d\theta$, and then subtract the differential of a cutoff function to make
the difference vanish there.  Average the construction over $\langle\tau,\sigma\rangle$.
The group acts transitively on the corners, so the four oriented boundary periods agree.

Adjust the primitive in disjoint boundary collars so that its Liouville vector field is outward
transverse, and attach the positive symplectizations of the boundary circles.  These adjustments
can be made equivariantly and do not change the primitive on the central annulus.
\end{proof}

Put
\[
 \Lambda_P=-\lambda_1-\lambda_2+\lambda_3.
\]
On a smooth torus-cover chart of the unperturbed correspondence, Smith's boundary coordinates are
\begin{equation}\label{eq:smith-boundary}
 \gamma_1=\gamma_2=\gamma_3=\gamma,
 \qquad
 \theta_1=\beta\sin\alpha,
 \qquad
 \theta_2=\theta-\beta\sin\alpha,
 \qquad
 \theta_3=\theta
\end{equation}
\cite[Eqs.~(3.1.4)--(3.1.5)]{Smith}.  At $\gamma=\pi/2$, the normalized local primitive therefore
satisfies $\Lambda_P|_{C_0}=0$.

\begin{proposition}[Exactness of the small perturbation]\label{prop:perturbed-exactness}
For a generic symmetric truncation $P^\delta$ and every sufficiently small $t>0$,
\[
 [\ell_t^*\Lambda_P]=0\quad\text{in }H^1(C_t^\delta;\R),
\]
where $\ell_t:C_t^\delta\iss(P^\delta)^-\times(P^\delta)^-\times P^\delta$ is the restriction
immersion.  The assertion includes the cycles created when Smith's singular seam is smoothed and
all corner-link cycles.
\end{proposition}

\begin{proof}
We first calculate on the full singular parameter fiber.  Smith's regular parameter space is a
two-fold cover of $C_0$ and decomposes as
\[
 \mathcal V_0^*\cong
 \bigl[S^2\times(S^1\setminus S^0)\bigr]
 \sqcup
 \bigl[T^2\times(S^1\setminus S^0)\bigr]
\]
\cite[Def.~3.6]{Smith}.  The first summand has no first real cohomology.  On the second, represent
the two torus generators at $\gamma=\pi/2$.  Equation~\eqref{eq:smith-boundary} gives angular
increments $(2\pi,-2\pi,0)$ for the $\beta$-circle and $(0,2\pi,2\pi)$ for the $\theta$-circle.
The normalized primitive vanishes on both generators.  Thus the period class vanishes on the
regular part.

It remains to include cycles through the singular seam.  The full zero set is
\[
 \mathcal V_0=A\cup H_0\cup H_\pi,
 \qquad
 A=\{\alpha=\pi/2\}\cong T^3,
 \qquad
 H_0\sqcup H_\pi=\{\gamma=0,\pi\}\cong S^0\times S^1\times S^2.
\]
Cut $A$ along the four subtori $\beta=\epsilon\pi$ and
$\theta-\beta=\epsilon\pi$, $\epsilon\in\{0,1\}$.  Mayer-Vietoris gives a spanning set for
first homology consisting of the angular directions on these subtori, four $\gamma$-circles, and
the two directions on every corner link.  The angular periods vanish by the preceding paragraph.

Orient the four pillowcase meridians positively and use
$(m_{0\pi},m_{\pi0},m_{\pi\pi})$ as a basis.  If
$\mathcal A_\delta=\int_{P^\delta}d\lambda$, direct substitution in the boundary formulas gives
\begin{equation}\label{eq:singular-gamma-table}
\begin{array}{c|c|c}
c&[-p_1(c)-p_2(c)+p_3(c)]&
 \displaystyle\int_c r_0^*\Lambda_P^{(00)}\\ \hline
A_0^\gamma&(0,0,0)&0\\
A_\pi^\gamma&(2,0,2)&\mathcal A_\delta\\
B_0^\gamma&(0,1,-1)&0\\
B_\pi^\gamma&(-2,-1,-1)&-\mathcal A_\delta.
\end{array}
\end{equation}
Here $\lambda^{(00)}$ is a primitive with meridian periods
$(-\mathcal A_\delta,0,0,0)$.  The symmetric primitive has period
$-\mathcal A_\delta/4$ on each meridian.  Pairing the difference with the middle column changes
the final column by $0,-\mathcal A_\delta,0,\mathcal A_\delta$, respectively, so all four
$\gamma$-periods vanish.

At a parameter corner $(g,h,k)\in\{0,1\}^3$, an oriented basis $(u,v)$ of its torus link has
meridian winding triples
\begin{equation}\label{eq:corner-link-table}
\begin{array}{c|c|c}
(g,h)&(w_1,w_2,w_3)(u)&(w_1,w_2,w_3)(v)\\ \hline
(0,0)&(-1,1,0)&(1,1,2)\\
(0,1)&(-1,-1,-2)&(-1,1,0)\\
(1,0)&(1,1,2)&(-1,1,0)\\
(1,1)&(1,-1,0)&(1,1,2).
\end{array}
\end{equation}
Every vector satisfies $w_3=w_1+w_2$.  Since the four end residues agree, the signed period of a
link circle is a common multiple of $-w_1-w_2+w_3$ and hence vanishes.  We have therefore proved
that $r_0^*\Lambda_P$ has zero period on every class in
$H_1(\mathcal V_0^\delta;\R)$.

For $t>0$, cut out the perturbation solid torus.  The perturbed correspondence is a geometric
composition of a fixed restriction immersion with the solid-torus arc
\[
 [e^{is}]\longmapsto[e^{is},e^{if_t(s)}],
\]
which is a Hamiltonian deformation \cite[Eq.~(6-4) and Sec.~6B]{CHK}.  If $v_t$ is the velocity of
the composed family, the Lagrangian condition gives
\begin{equation}\label{eq:hamiltonian-flux}
 \ell_t^*\iota_{v_t}(-\omega_1-\omega_2+\omega_3)
 =-d(H_t\circ\operatorname{pr}),
\end{equation}
with the convention $\iota_{X_{H_t}}\omega=-dH_t$.  Hence
$\ell^*\Lambda_P+G\,dt$ is closed on the smooth trace of the family.

Let $q_t:\mathcal V_t^\delta\to C_t^\delta$ be Smith's two-fold parameter cover.  The closure of
its trace over $0\le t\le\epsilon$ is a compact subanalytic pair with the zero fiber, so it admits
a compatible triangulation.  After decreasing $\epsilon$, a regular neighborhood of the zero
fiber contains the trace and retracts onto that fiber.  Thus every one-cycle in a positive fiber
is joined by a singular two-chain to a cycle in the full zero set; vanishing cycles are included.
The Hamiltonian graph and all boundary maps extend analytically to $t=0$, so the trace form extends
continuously and is smooth on each stratum of the triangulation.  Stokes' theorem applied stratum
by stratum gives
\[
 \int_{\widetilde a_t}q_t^*\ell_t^*\Lambda_P
 =\int_{a_0}r_0^*\Lambda_P=0.
\]
Pullback on first real cohomology under a finite cover is injective.  This proves the proposition.
\end{proof}

The rational arcs are contractible and can be normalized equivariantly, so
$A_{1/3}\times A_{1/7}$ is exact.  Primitives add under transverse geometric composition.
Proposition~\ref{prop:perturbed-exactness} therefore gives a primitive on every component of the
direct fiber product \eqref{eq:composition}.  The output, however, passes through two deleted
corners in each lifted period.  Its two open arcs initially have independent additive constants;
the next proposition fixes their relative normalization.

\begin{proposition}[Half-period normalization]\label{prop:half-period}
On the parameter cover of the direct composition, all-corner points on the main lift are
\[
 z_j=(\gamma,\theta,\alpha,\beta)
     =(21j\pi,10j\pi,\pi/2,7j\pi),
 \qquad j\in\Z.
\]
Removing two consecutive corner points cuts a lifted period into two open arcs.  Define
\[
 T(\gamma,\theta,\alpha,\beta)
  =(\gamma+21\pi,\theta+10\pi,\alpha,\beta+7\pi).
\]
Modulo the parameter period lattice, $T$ is an involution exchanging the two arcs.  With
$\rho=\tau\sigma$, the boundary maps satisfy
\begin{equation}\label{eq:half-period-boundary}
 r_1T=\rho r_1,
 \qquad r_2T=\rho r_2,
 \qquad r_3T=\sigma r_3.
\end{equation}
The inherited primitive $F$ may be normalized so that $F\circ T=F$.

Let $I_0$ be the open arc $0<\gamma<21\pi$, let $I_1=T(I_0)$, and define
\[
 \operatorname{can}(z)=z\quad(z\in I_0),
 \qquad
 \operatorname{can}(z)=T^{-1}z\quad(z\in I_1).
\]
For every root $x=(z_-,z_+)$, including a root whose preimages lie on different open arcs,
\begin{equation}\label{eq:canonical-action}
 \mathcal A_t(x)
 =\int_{\operatorname{can}(z_+)}^{\operatorname{can}(z_-)}r_3^*\lambda,
\end{equation}
where the integral follows the oriented subarc of $I_0$.
\end{proposition}

\begin{proof}
At a parameter corner $(g,h,k)$, the three boundary characters are
$(g,k)$, $(g,k+h)$, and $(g,h)$.  For odd $q$, a corner $(g,r)$ lies on $A_{1/q}$ precisely when
$r=g$.  The two input equations therefore give $k=g$ and $h=0$, which yields exactly the displayed
corner lifts.  Their consecutive spacing shows that each full period contains two corner cuts.

The square of $T$ is the parameter period $(42\pi,20\pi,0,14\pi)$.  Direct substitution in the
defining traceless equation shows that $T$ preserves the parameter space, and substitution in the
boundary words gives \eqref{eq:half-period-boundary}.  The map $\rho$ preserves both odd-slope
arcs.  Average the primitive on each rational arc under $\rho$ and the primitive on $C_t$ under
$T$.  The normalized Liouville form is invariant under $\tau$ and $\sigma$, so the averages have
the same exterior derivatives and the composition primitive satisfies $F\circ T=F$.

On either open component, $dF=r_3^*\lambda$.  Replace both endpoints by their representatives in
$I_0$ and integrate from the positive endpoint to the negative endpoint.  The interior of $I_0$
contains no all-corner point, so no residue detour is involved and
\eqref{eq:canonical-action} follows.
\end{proof}

\section{The six folds and the separated graph sheets}\label{sec:geometry}

The partial composition
\[
 F_t=A_{1/3}\times_{P_1}C_t\iss P_2^-\times P_3
\]
is used here only to describe the finite geometry of the direct output.  No associativity or
curved-composition theorem is invoked.

\begin{proposition}[Six-circle and fold table]\label{prop:six-fold-table}
The singular partial composition has six stabilizer circles.  Their $P_2$ seam values, the
corresponding $A_{1/3}$ angle $\theta_1$, and the image interval $[\ell_i,h_i]$ in the third
pillowcase are
\[
\begin{array}{c|c|c|c}
i&(\gamma_2,\theta_2)&\theta_1&[\ell_i,h_i]\\ \hline
1&(0,2\pi/7)&2\pi/3&[8\pi/21,20\pi/21]\\
2&(0,4\pi/7)&2\pi/3&[2\pi/21,16\pi/21]\\
3&(0,6\pi/7)&2\pi/3&[4\pi/21,10\pi/21]\\
4&(\pi,\pi/7)&\pi/3&[4\pi/21,10\pi/21]\\
5&(\pi,3\pi/7)&\pi/3&[2\pi/21,16\pi/21]\\
6&(\pi,5\pi/7)&\pi/3&[8\pi/21,20\pi/21].
\end{array}
\]
If $\psi$ parametrizes a stabilizer circle, then
\begin{equation}\label{eq:circle-fold}
 \cos\theta_3(\psi)
 =\cos\theta_1\cos\theta_2
  +\sin\theta_1\sin\theta_2\cos\psi.
\end{equation}
Both endpoints of every interval are nondegenerate folds.  The pillowcase symmetry pairs rows
$(1,6)$, $(2,5)$, and $(3,4)$, including their fold coefficients.
\end{proposition}

\begin{proof}
A lift of $A_{1/q}$ satisfies $q\theta-\gamma\in2\pi\Z$.  Its noncorner seam values are
\[
 A_{1/3}:\ (0,2\pi/3),\ (\pi,\pi/3)
\]
and
\[
 A_{1/7}:\ (0,2\pi k/7),\ (\pi,(2k-1)\pi/7),
 \qquad k=1,2,3.
\]
Pairing the unique $A_{1/3}$ value on each seam with the three $A_{1/7}$ values gives the six
circles.  Choose the common seam axis as polar axis.  The Euclidean dot product of the two
remaining axes gives \eqref{eq:circle-fold}, and hence
\[
 [\ell_i,h_i]
 =\bigl[|\theta_1-\theta_2|,
         \pi-|\pi-(\theta_1+\theta_2)|\bigr].
\]
At the two endpoints,
\begin{align*}
 \theta_3(\psi)&=\ell_i+
 \frac{\sin\theta_1\sin\theta_2}{2\sin\ell_i}\psi^2+O(\psi^4),\\
 \theta_3(\pi+u)&=h_i-
 \frac{\sin\theta_1\sin\theta_2}{2\sin h_i}u^2+O(u^4).
\end{align*}
All coefficients are positive and finite.  The symmetry assertion follows by replacing the two
input angles by their complements.
\end{proof}

\begin{proposition}[Hyperbolic smoothing at the twelve endpoints]
\label{prop:hyperbolic-normal-form}
Use the group-presentation convention
\[
 p=\exp\!\bigl(t\,\operatorname{Im}(ac^{-1}x)\bigr),
 \qquad
 \Phi_t=\operatorname{Re}(xp^{-1}a^{-1}pb),
\]
with
\[
 a=\mathbf i,
 \quad b=e^{\gamma\mathbf k}\mathbf i,
 \quad c=e^{\theta\mathbf k}\mathbf i,
 \quad
 x=\sin\alpha\cos\beta\,\mathbf i
   +\sin\alpha\sin\beta\,\mathbf j+\cos\alpha\,\mathbf k.
\]
Near the lower and upper endpoints of the $i$th circle there are smooth coordinates $(u,v)$ in
which
\[
 \{\Phi_t=0\}=\{uv=\tau_i^\pm(t)\},
 \qquad
 \tau_i^\pm(0)=0,
 \qquad
 (\tau_i^\pm)'(0)=\mp2\sin\theta_1\sin\theta_2.
\]
Thus all twelve endpoints have nondegenerate hyperbolic smoothing germs, with opposite sectors at
the two ends of each circle and the same three symmetry pairs as above.
\end{proposition}

\begin{proof}
At $t=0$ the defining equation is
\[
 \Phi_0(\gamma,\theta,\alpha,\beta)=\sin\gamma\cos\alpha.
\]
Write $X=\gamma-\gamma_i$ and $Y=\alpha-\pi/2$, where
$\gamma_i\in\{0,\pi\}$ and $\epsilon_i=\cos\gamma_i$.  Quaternion multiplication gives
\[
 \Phi_t
 =\epsilon_i\bigl(-XY-2t\sin\beta\sin(\theta-\beta)\bigr)
  +O\!\left(t(|X|+|Y|)+t^2+(|X|+|Y|)^3\right).
\]
The Hessian in $(X,Y)$ is nonsingular with index one.  The parametric Morse lemma therefore gives
$uv=\tau(t)$ and
\[
 \tau'(0)=-2\sin\beta\sin(\theta-\beta).
\]
At the lower and upper lifts the products are respectively
$-\sin\theta_1\sin\theta_2$ and $\sin\theta_1\sin\theta_2$, proving the formulas.
\end{proof}

\begin{proposition}[Off-seam graph sheets]\label{prop:graph-sheets}
Away from neighborhoods of the six seam values, every branch of $F_0$ is locally the graph of the
symplectic shear
\[
 \phi_c(\gamma,\eta)=(\gamma,\eta+\gamma/3+c).
\]
On a torus-cover chart, the $21$ local lifts of the direct composition are parallel and pairwise
disjoint.  Both statements persist on compact off-seam subsets for all sufficiently small $t>0$.
\end{proposition}

\begin{proof}
On the unperturbed main stratum, the three boundary coordinates are
\[
 p_1=(\gamma,\beta),
 \qquad p_2=(\gamma,\theta-\beta),
 \qquad p_3=(\gamma,\theta).
\]
The first rational arc gives $\beta=\gamma/3+c$.  With $\eta=\theta-\beta$, this is the displayed
shear, and
$d\gamma\wedge d(\eta+\gamma/3+c)=d\gamma\wedge d\eta$.
Smith's cross-resolution theorem confines every failure of this graph description to the six seam
neighborhoods \cite[Thm.~4.22]{Smith}.

Choose lifts
\[
 \beta=\frac{\gamma+2\pi k}{3},
 \qquad
 \eta=\frac{\gamma+2\pi\ell}{7},
 \qquad (k,\ell)\in\Z/3\times\Z/7.
\]
Then
\[
 \theta=\frac{10}{21}\gamma+
 \frac{2\pi(7k+3\ell)}{21}.
\]
The homomorphism $(k,\ell)\mapsto7k+3\ell$ from
$\Z/3\times\Z/7$ to $\Z/21$ is bijective, so the $21$ lifts have distinct intercepts.  A compact
off-seam set has positive minimum separation, which persists under the $C^1$ small-$t$
deformation.
\end{proof}

\section{The exact root-action table}\label{sec:actions}

Use scaled coordinates on the canonical arc $I_0$:
\[
 s=\frac{\gamma}{\pi},
 \qquad
 n=\frac{21\theta}{\pi}.
\]
A vertical torus period changes $n$ by $42$, and the arc runs from $(0,0)$ to $(21,210)$.

\begin{lemma}[Singular first-arc word]\label{lem:path-word}
Between consecutive integral values of $s$, the singular first arc has slope
$d\theta/d\gamma=10/21$.  Its twelve signed vertical jumps are
\[
\begin{array}{c|r|r|c@{\qquad}c|r|r|c}
s&i&\Delta n&n\text{-interval}&s&i&\Delta n&n\text{-interval}\\ \hline
1&4&-6 &[4,10]   &11&5& 14 &[110,124]\\
4&1&-12&[22,34]  &12&1& 12 &[134,146]\\
5&4& 6 &[32,38]  &14&2&-14 &[152,166]\\
7&5&-14&[44,58]  &16&3& 6  &[172,178]\\
9&6& 12&[64,76]  &17&6&-12 &[176,188]\\
10&2&14&[86,100] &20&3&-6  &[200,206].
\end{array}
\]
At the other interior integral seams, the diagonal sheet has the eight points
\[
 (2,14),(3,24),(6,48),(8,54),
 (13,156),(15,162),(18,186),(19,196).
\]
\end{lemma}

\begin{proof}
Proposition~\ref{prop:graph-sheets} gives
$\theta=\gamma/3+\gamma/7+c$, so increasing $s$ by one increases $n$ by ten.  At an integral
seam, substitute the six intervals of Proposition~\ref{prop:six-fold-table} and choose the lift
whose incoming endpoint agrees with the preceding diagonal segment.  Continuing through the
outgoing endpoint gives the twelve jumps.  The recurrence
\[
 n_j^-=n_{j-1}^++10,
 \qquad
 n_j^+=n_j^-+\Delta n_j
\]
gives all displayed intervals and point values.  The jumps sum to zero, so the endpoint is
$(21,210)$.
\end{proof}

\begin{lemma}[Algebraic normal roots]\label{lem:normal-roots}
Put $x=\cos\theta$ and $u=2\cos(\pi/7)$, so
\[
 u^3-u^2-2u+1=0.
\]
On row $i$ of Proposition~\ref{prop:six-fold-table}, the small-$t$ normal displacement is
\begin{equation}\label{eq:normal-displacement}
 \gamma=j\pi+tG_i(\theta)+O(t^2),
 \qquad
 G_i(\theta)=\frac{2(a_i-\cos\theta)}{\cos\alpha_i(\theta)},
\end{equation}
where
\begin{equation}\label{eq:alpha-square}
 \cos^2\alpha_i(\theta)
 =\frac{b_i^2-(\cos\theta-a_i)^2}{\sin^2\theta},
 \qquad
 \operatorname{sgn}(\cos\alpha_i)=
 \begin{cases}+1,&i=1,2,3,\\-1,&i=4,5,6,
 \end{cases}
\end{equation}
and
\[
\begin{gathered}
 a_1=a_6=\frac{2-u^2}{4},
 \qquad
 a_2=a_5=\frac{u^2-u-1}{4},
 \qquad
 a_3=a_4=\frac u4,\\
 b_i^2=\frac34(1-4a_i^2).
\end{gathered}
\]
The physical simple solutions of $G_i=G_j$ on $0<\theta<\pi$ are exactly
\[
\begin{array}{c|c|c|c}
\text{root}&\text{row pairs }\{i,j\}&x\text{-box}
 &21\theta/\pi\text{-box}\\ \hline
\mathsf A&\{1,4\},\{3,6\}&[0.178447,0.178449]&[9.300,9.301]\\
\mathsf B&\{1,5\},\{2,6\}&[-0.123491,-0.123489]&[11.327,11.328]\\
\mathsf C&\{1,6\}&[-0.311746,-0.311743]&[12.619,12.620]\\
\mathsf D&\{2,3\},\{4,5\}&[0.722519,0.722522]&[5.102,5.103]\\
\mathsf E&\{2,4\},\{3,5\}&[0.346009,0.346012]&[8.138,8.139]\\
\mathsf F&\{2,5\}&[0.111259,0.111262]&[9.754,9.755]\\
\mathsf G&\{3,4\}&[0.450483,0.450486]&[7.376,7.377].
\end{array}
\]
There is no physical root for the row pairs
$\{1,2\},\{1,3\},\{4,6\},\{5,6\}$.
\end{lemma}

\begin{proof}
At $t=0$, substitution of the seam rational-arc equations into the quaternionic boundary words
gives \eqref{eq:alpha-square}.  Differentiating the defining equation in the normal $\gamma$
direction and in $t$ gives
\[
 \partial_\gamma\Phi_0=(-1)^j\cos\alpha_i,
 \qquad
 \partial_t\Phi_0=-2(-1)^j(a_i-\cos\theta),
\]
which proves \eqref{eq:normal-displacement}.  The formulas for $a_i$ and $b_i$ follow from the
minimal polynomial of $u$.

For pairs of distinct symmetry types, squaring $G_i=G_j$ and clearing the positive fold
denominators gives the three quadratics $P_{rs}(x)=A_{rs}x^2+B_{rs}x+C_{rs}$ with
\begin{align*}
 (A_{12},B_{12},C_{12})
 &=\left(u^2-2u-1,
   -\frac52u^2+\frac32u+\frac52,
   \frac14u^2-\frac12u-\frac14\right),\\
 (A_{13},B_{13},C_{13})
 &=\left(u^2-\frac12u-\frac32,
   -\frac34u^2-\frac12u+2,
   \frac14u^2-\frac18u-\frac38\right),\\
 (A_{23},B_{23},C_{23})
 &=\left(u^2+u-2,
   u^2-\frac52u-2,
   \frac14u^2+\frac14u-\frac12\right).
\end{align*}
The rational interval
\[
 1.801937735<u<1.801937737
\]
contains $2\cos(\pi/7)$ and one simple zero of $u^3-u^2-2u+1$.  Rational interval evaluation
gives one physical root of $P_{12}$, one of $P_{13}$, and both roots of $P_{23}$ in the displayed
boxes.  Each polynomial has opposite strict signs at the endpoints of its box and a derivative of
one strict sign there.  Vieta's formula places the unused root of $P_{12}$ below $-1$ and that of
$P_{13}$ above $1$.  The signs in \eqref{eq:alpha-square} remove the extraneous squared solutions.
For equal symmetry types the two functions agree only at $x=a_i$, producing
$\mathsf C,\mathsf F,\mathsf G$, and their derivatives are nonzero.

The angle boxes are also rationally certified.  Machin's identity
\[
 \frac\pi4=4\arctan\frac15-\arctan\frac1{239}
\]
and alternating sums give a rational $\pi$-interval of width less than
$1.84\times10^{-35}$.  Rational Taylor bounds for cosine then place every angle strictly inside
its displayed box and certify the derivative signs.
\end{proof}

\begin{proposition}[Complete limiting action table]\label{prop:action-table}
The path of Lemma~\ref{lem:path-word} has $58$ seam-overlap candidates.  Exactly $50$ persist as
transverse degree-one roots for all sufficiently small $t>0$: $28$ are diagonal-circle roots and
$22$ are the simple circle-circle roots of Lemma~\ref{lem:normal-roots}.  The eight excluded seam
pairs are
\[
 \{1,17\},\{4,14\},\{4,20\},\{5,9\},
 \{7,17\},\{9,11\},\{10,12\},\{12,16\}.
\]

The following table gives one representative of each half-period pair.  Here $d\in42\Z$ is the
offset in $n$, $r$ is a normal root or $\mathrm{pt}$ for a diagonal-circle root, $D$ is the short
deck class defined above, and $o$ is the degree-one ordering.  The columns $a,w,\kappa$ satisfy
\begin{equation}\label{eq:area-residue}
 \frac1{\pi^2}\int_j^k\lambda=\frac\kappa2(a-w).
\end{equation}
\begingroup
\scriptsize
\setlength{\tabcolsep}{3.2pt}
\[
\begin{array}{c|r|c|c|c|r|r|c|r}
(j,k)\sim(j',k')&d&r&D&o&a&w&\kappa&\mathcal A/\pi^2\\ \hline
(1,6)\sim(15,20)&42&\mathrm{pt}&(8,4)&1\to6&80/21&2&1/2&-19/42\\
(1,7)\sim(14,20)&42&\mathsf D&(3,1)&7\to1&-8/7&-2&1&3/7\\
(1,10)\sim(11,20)&84&\mathsf E&(6,3)&1\to10&124/7&14&1/2&-13/14\\
(1,12)\sim(9,20)&126&\mathsf A&(5,2)&1\to12&288/7&38&1/2&-11/14\\
(1,16)\sim(5,20)&168&\mathsf G&(3,1)&1\to16&36&34&1/2&-1/2\\
(2,7)\sim(14,19)&42&\mathrm{pt}&(8,4)&7\to2&24/7&2&1/2&5/14\\
(2,10)\sim(11,19)&84&\mathrm{pt}&(4,2)&2\to10&96/7&12&1&-6/7\\
(2,12)\sim(9,19)&126&\mathrm{pt}&(5,3)&2\to12&192/7&26&1&-5/7\\
(2,17)\sim(4,19)&168&\mathrm{pt}&(3,1)&17\to2&244/7&34&1/2&3/14\\
(3,4)\sim(17,18)&0&\mathrm{pt}&(10,5)&4\to3&-116/21&-6&1/2&5/42\\
(3,9)\sim(12,18)&42&\mathrm{pt}&(3,1)&3\to9&-16/21&-2&1&-13/21\\
(4,5)\sim(16,17)&0&\mathsf A&(10,5)&4\to5&-36/7&-6&1/2&-3/14\\
(4,9)\sim(12,17)&42&\mathsf C&(8,4)&4\to9&8&6&1/2&-1/2\\
(4,11)\sim(10,17)&84&\mathsf B&(7,3)&4\to11&284/7&38&1/2&-9/14\\
(4,13)\sim(8,17)&126&\mathrm{pt}&(6,2)&4\to13&1496/21&70&1/2&-13/42\\
(5,11)\sim(10,16)&84&\mathsf D&(3,2)&5\to11&188/7&26&1&-3/7\\
(5,14)\sim(7,16)&126&\mathsf E&(6,2)&14\to5&488/7&70&1/2&-1/14\\
(5,15)\sim(6,16)&126&\mathrm{pt}&(5,3)&15\to5&544/21&26&1&-1/21\\
(6,7)\sim(14,15)&0&\mathrm{pt}&(10,5)&7\to6&-212/21&-10&1/2&-1/42\\
(6,10)\sim(11,15)&42&\mathrm{pt}&(2,1)&6\to10&188/21&8&1&-10/21\\
(7,8)\sim(13,14)&0&\mathrm{pt}&(10,5)&7\to8&-28/3&-10&1/2&-1/6\\
(7,10)\sim(11,14)&42&\mathsf F&(9,4)&7\to10&32&30&1/2&-1/2\\
(7,12)\sim(9,14)&84&\mathsf B&(8,3)&7\to12&500/7&70&1/2&-5/14\\
(8,10)\sim(11,13)&42&\mathrm{pt}&(1,1)&8\to10&68/3&22&1&-1/3\\
(8,12)\sim(9,13)&84&\mathrm{pt}&(2,2)&8\to12&932/21&44&1&-4/21.
\end{array}
\]
\endgroup
Thus $42$ roots have negative limiting action and $8$ have positive limiting action.  Every
limiting action is nonzero, and the minimum absolute value is $\pi^2/42$.
\end{proposition}

\begin{proof}
For seams $j<k$, canonicalized output points agree precisely when their $n$-intervals overlap
after a shift $d\in42\Z$.  Intersecting the twenty intervals in Lemma~\ref{lem:path-word} gives
$58$ nonempty triples $(j,k,d)$.  Twenty-eight contain one diagonal point strictly inside a
vertical interval.  The other thirty are circle-circle overlaps.  Substitution of the certified
normal roots from Lemma~\ref{lem:normal-roots} leaves one simple root in $22$ and none in the eight
listed pairs.  The strict angle boxes lie inside their overlap intervals.

At a diagonal-circle root, the limiting directions in scaled coordinates are
$(1,10/21)$ and $(0,\Delta n_j)$.  At a circle-circle root in rows $i,i'$, the determinant has
leading sign
\[
 \operatorname{sgn}(\Delta n_j\Delta n_k)
 \operatorname{sgn}\bigl(G_i'(\theta)-G_{i'}'(\theta)\bigr).
\]
The rational derivative certificates therefore fix the degree-one ordering and prove
transversality.  Proposition~\ref{prop:graph-sheets} excludes additional off-seam roots, while
Proposition~\ref{prop:hyperbolic-normal-form} gives an embedded local resolution at each fold.

For an overlap $(j,k,d)$ choose a common height $h$ and put
\[
 p=(j,h/21),
 \qquad
 q=(k,(h+d)/21)
\]
in $(\gamma/\pi,\theta/\pi)$ coordinates.  Follow the singular path from $p$ to $q$.  If $k-j$ is
even, its endpoints differ by an ordinary torus deck transformation and $\kappa=1$.  If $k-j$ is
odd, concatenate the path with its endpoint translate and divide the resulting integral by two;
then $\kappa=1/2$.

Concatenate this ordinary loop with its pillowcase-involuted image and close the remaining ends by
mutually cancelling auxiliary arcs.  Stokes' theorem and the residue $-\pi^2$ at each lattice point
give
\[
 \frac1{\pi^2}\int_p^q\lambda
 =\frac\kappa2\left(
   \operatorname{Area}-\sum_{z\in\Z^2}\operatorname{wind}_z\right),
\]
which is \eqref{eq:area-residue}.  Shoelace evaluation on the vertices of
Lemma~\ref{lem:path-word} gives the printed $a$ and $w$.  The common height can move inside its
overlap interval without crossing a lattice point, so the integral is constant and may be
evaluated at a rational midpoint even when the actual normal root is algebraic.

For example, $(1,7)$ has $(a,w,\kappa)=(-8/7,-2,1)$, so the $7\to1$ action is
$3\pi^2/7$.  The pair $(3,4)$ has
$(a,w,\kappa)=(-116/21,-6,1/2)$, so the $4\to3$ action is $5\pi^2/42$.  The half-period map sends
\[
 (j,k,d)\longmapsto(21-k,21-j,d)
\]
and preserves the normalized action, ordering, and short deck class.  Each row therefore has
multiplicity two.  Reading the last column proves the sign census and the minimum gap.
\end{proof}

\begin{corollary}[Small-perturbation persistence]\label{cor:persistence}
There is $t_0>0$ such that, for every $0<t<t_0$, the curve $L_t$ has exactly the $50$ roots in
Proposition~\ref{prop:action-table}, and all have the signs shown there.
\end{corollary}

\begin{proof}
Transversality continues all $50$ roots uniquely.  Compact separation of the graph sheets and the
fold neighborhoods excludes new roots.  Proposition~\ref{prop:half-period} gives primitives that
vary continuously with $t$, and the canonical paths remain in the compact corner-free arc $I_0$.
Thus every action converges to its table value.  Since the finite minimum absolute limiting value
is $\pi^2/42$, every sign is constant after decreasing $t_0$.
\end{proof}

\begin{proof}[Proof of Theorem~\ref{thm:main}]
Exactness is Proposition~\ref{prop:perturbed-exactness}, and the relative primitive is fixed by
Proposition~\ref{prop:half-period}.  Proposition~\ref{prop:action-table} proves the complete root
and limiting-action census.  Corollary~\ref{cor:persistence} proves persistence for small positive
$t$.
\end{proof}

\section{Energy interpretation and scope}\label{sec:scope}

The sign table has an elementary consequence for any exact holomorphic configuration that may be
defined using $L_t$.  The statement does not assume that a compactified figure-eight operation
exists.

\begin{lemma}[Root-action identity]\label{lem:energy-action}
Let a finite tree of exact pseudoholomorphic or Morse components have no external inputs and one
external root $x$ on $L_t$.  Suppose its seam and boundary lifts are fixed and every component has
the usual Stokes energy identity.  Then
\[
 E_{\mathrm{tot}}=\mathcal A_t(x).
\]
Consequently, a nonconstant such configuration can be rooted only at a positive-action root.
\end{lemma}

\begin{proof}
Apply Stokes' theorem to each component.  At every internal flag the same primitive difference
appears with opposite signs on the two adjacent components.  The action drops along Morse edges
cancel in the same way.  With no external inputs, only the external output remains.  Compatible
almost-complex and Morse data make each energy summand nonnegative, and a nonconstant component
makes the total positive.
\end{proof}

\begin{corollary}\label{cor:eight-candidates}
For $0<t<t_0$, every nonconstant exact zero-input configuration satisfying the hypotheses of
Lemma~\ref{lem:energy-action} is rooted at one of the eight positive roots in
Theorem~\ref{thm:main}.
\end{corollary}

\begin{remark}[No counting conclusion]\label{rem:no-count}
Corollary~\ref{cor:eight-candidates} is an action obstruction, not an existence or parity theorem.
It does not assert that any of the eight positive roots supports a holomorphic configuration, that
only one vertex occurs, or that a sum of outputs defines a bounding cochain.  Those conclusions
would require compactification, transversality, gluing, local mod-two counts, and continuation
data that are not constructed here.
\end{remark}

\begin{remark}[Exact certificate]\label{rem:certificate}
The command
\[
 \texttt{python3 pillowcase/q7\_exact\_actions.py}
\]
reconstructs the path word, verifies the rational root enclosures and derivative signs, enumerates
all $58$ overlaps, and prints the $25$-row area-residue table using exact rational arithmetic.
The standard-library program is an independent reproducibility certificate for the finite
calculations, not a numerical premise of the proof.
\end{remark}

Theorem~\ref{thm:main} therefore supplies a complete finite action filtration for this composition.
It leaves the holomorphic and categorical questions where they belong: outside the claimed result.

\raggedright
\begin{thebibliography}{CHKK22}

\bibitem[BW18]{BottmanWehrheim}
N.~Bottman and K.~Wehrheim,
\emph{Gromov compactness for squiggly strip shrinking in pseudoholomorphic quilts},
Selecta Math. (N.S.) \textbf{24} (2018), 3381--3443. arXiv:1503.03486.

\bibitem[CHKK22]{CHKK}
G.~Cazassus, C.~Herald, P.~Kirk, and A.~Kotelskiy,
\emph{The correspondence induced on the pillowcase by the earring tangle},
J. Topol. \textbf{15} (2022), 2472--2543. arXiv:2010.04320.

\bibitem[CHK22]{CHK}
G.~Cazassus, C.~Herald, and P.~Kirk,
\emph{Tangles, relative character varieties, and holonomy perturbed traceless flat moduli spaces},
Open Book Ser. \textbf{5} (2022), 95--121. arXiv:2102.00901.

\bibitem[HK18]{HeraldKirk}
C.~Herald and P.~Kirk,
\emph{Holonomy perturbations in a cylinder, and regularity for traceless
$\mathrm{SU}(2)$ character varieties of tangles},
Geom. Topol. \textbf{22} (2018), 129--216. arXiv:1511.00308.

\bibitem[Smi24]{Smith}
K.~Smith,
\emph{Perturbed traceless $\mathrm{SU}(2)$ character varieties of tangle sums},
arXiv:2412.06066.

\bibitem[Wue26]{WuebbenII}
B.~J.~Wuebben,
\emph{The instanton homology of the $(-2,3,q)$ pretzel knots and Maurer--Cartan deformations in
the two-arc algebra of the pillowcase},
arXiv:2607.26096.

\end{thebibliography}

\bigskip
\noindent{\sc Bernd J. Wuebben, New York, NY,} \texttt{wuebben@gmail.com}

\end{document}
